%% kapa: exam=1 round=27 year=2016 period=1 form=A subject=civil_law subject_ko=민법 seq=34 number=34
%% source: Q-Net 감정평가사 기출문제 (cst00309) artlSeq=5169402 file=extracted/1008/5169402/제27회 감정평가사 1차 문제/제27회 감정평가사 1교시문제(A형).pdf
%% conversion: kordoc 4.24.51 -> tools/gen1.py (auto); review=proofread-sonnet; flags=none
\begin{question}{34}
근저당권에 관한 설명으로 옳은 것은? (다툼이 있으면 판례에 따름)
\begin{choices}
  \item 근저당권의 피담보채무가 확정되기 이전에는 채무자를 변경할 수 없다.
  \item 근저당권의 확정 전에 발생한 원본채권으로부터 그 확정 후에 발생하는 이자는 채권최고액의 범위 내에서 여전히 담보된다.
  \item 선순위근저당권자가 경매를 신청하는 경우, 후순위근저당권의 피담보채권의 확정시기는 경매개시결정시이다.
  \item 근저당권의 존속 중에 피담보채권이나 기본계약과 분리하여 근저당권만을 양도 할 수도 있다.
  \item 채권의 총액이 채권최고액을 초과하는 경우, 채무자 겸 근저당권설정자는 근저당권의 확정 전이라도 채권최고액을 변제하고 근저당권의 말소를 청구할 수 있다.
\end{choices}
\end{question}
\answer{2}
